% Generated draft: requires mathematical and source review.
\documentclass[11pt]{article}
\usepackage[margin=25mm]{geometry}
\usepackage{fontspec}
\setmainfont{Latin Modern Roman}
\usepackage{amsmath,amssymb,mathtools,unicode-math}
\setmathfont{Latin Modern Math}
\usepackage{enumitem,needspace,xurl,hyperref}
\setlength{\parindent}{0pt}
\setlength{\parskip}{5pt}
\setlength{\emergencystretch}{3em}
\newcommand{\N}{\mathbb{N}}
\newcommand{\Z}{\mathbb{Z}}
\newcommand{\Q}{\mathbb{Q}}
\newcommand{\R}{\mathbb{R}}
\newcommand{\C}{\mathbb{C}}
\newcommand{\F}{\mathbb{F}}
\newcommand{\A}{\mathbb{A}}
\newcommand{\PP}{\mathbb{P}}
\newcommand{\E}{\mathbb{E}}
\newcommand{\eps}{\varepsilon}
\newcommand{\dd}{\,\mathrm{d}}
\begin{document}
\section*{UnsolvedMath ID: 20002669}
\textbf{Source problem code:} catalogue entry
\subsection{Definitions and mathematical statement}
Let \((M,\tau)\) be a finite von Neumann algebra with faithful normal tracial state, generated by a finite self-adjoint tuple \(x_1,\ldots,x_m\). In the standard bounded-generator interpretation, the free difference quotients satisfy \(\partial_jx_i=\delta_{ij}1\otimes1\). Finite free Fisher information means that conjugate variables \(\xi_j\in L^2(M,\tau)\) exist with
\[
 \tau(\xi_jP)=(\tau\otimes\tau)(\partial_jP),\qquad
 \Phi^*(x_1,\ldots,x_m)=\sum_j\|\xi_j\|_2^2<\infty.
\]
Use the closed free-difference-quotient quadratic form, with its required domains, to interpret
\[
 \Delta=\sum_j\partial_j^*\partial_j\geq0,
 \qquad\phi_t=e^{-t\Delta}.
\]
The source treats \(\phi_t\) as a completely positive heat deformation; the form/Markov construction must be checked in the chosen setting.

\textbf{Source question.} Can this semigroup converge uniformly to the identity on the operator-norm unit ball of \(M\) as \(t\downarrow0\)? The aid identifies the intended output norm as \(L^2\), hence the target is
\[
 H(t)=\sup_{\|a\|_\infty\leq1}\|\phi_t(a)-a\|_2\longrightarrow0.
\]
This is stronger than pointwise \(L^2\) continuity and differs from operator-norm convergence. The source adds a proposed non-weak-rigidity consequence for a nonhyperfinite generated algebra; that consequence requires the appropriate rigidity theorem, not merely a change of norm convention.

\textbf{Reported general spectral criterion.} For any positive self-adjoint \(\Delta\) on \(L^2(M)\), let
\[
 Q(R)=\sup_{\|a\|_\infty\leq1}
       \|\mathbf1_{(R,\infty)}(\Delta)a\|_2,
\qquad
 V(t)=\sup_{\|a\|_\infty\leq1}
       \|(I+t\Delta)^{-1}a-a\|_2.
\]
The aid reports equivalence of \(H(t)\to0\), \(V(t)\to0\) as \(t\downarrow0\), and \(Q(R)\to0\) as \(R\to\infty\). Quantitative sufficient comparisons, using \(\|a\|_2\leq1\), are
\[
 H(t)\leq tR+Q(R),\qquad
 Q(R)\leq\frac{H(t)}{1-e^{-tR}},
\]
\[
 V(t)\leq tR+Q(R),\qquad
 Q(R)\leq\frac{1+tR}{tR}V(t)
 \quad(t,R>0).
\]
Thus it is uniform high-spectral-tail control of the whole operator unit ball that is at issue; unboundedness of \(\Delta\) alone does not answer the source question.

\textbf{Reported compact-resolvent obstruction.} If \(M\) is diffuse and \(\Delta\) has compact resolvent, the aid reports
\[
 H(t)\geq1\qquad(t>0).
\]
A diffuse finite algebra contains a Haar unitary \(u\), meaning \(\tau(u^k)=0\) for every nonzero integer \(k\). Its positive powers are an orthonormal \(L^2\) sequence, with \(\|u^k\|_\infty=\|u^k\|_2=1\). Compactness of the positive-time heat operator supplies the reported witnesses
\[
 \|e^{-t\Delta}u^k\|_2\to0,
 \qquad\|e^{-t\Delta}u^k-u^k\|_2\to1.
\]
This is a sufficient obstruction under compact resolvent; no compactness is inferred for every finite-Fisher tuple.

For a finite standard freely independent semicircular family, the aid reports that the difference-quotient generator is the free number operator, with eigenvalue word length and finite-dimensional eigenspaces. Its resolvent is compact, and the generated algebra is diffuse, so the semigroup is nonuniform in the stated sense, including the one-variable case.

The aid also cites Peterson's \(L^2\)-rigidity framework and Dabrowski--Ioana criteria for non-\(L^2\)-rigidity/non-weak-rigidity in broad finite-Fisher regimes, but supplies neither full hypotheses nor citations. In the usual terminology weak rigidity concerns a diffuse subalgebra with a relative rigidity/property-\(T\) condition, while hyperfiniteness concerns finite-dimensional approximation; the exact theorem connecting these notions is not reconstructed. The fully general finite-Fisher formulation is explicitly not claimed solved.
\subsection{Short English statement}
The question concerns uniform \(L^2\) approximation on the operator-norm unit ball as time tends to zero. The aid characterizes this by uniform spectral-tail decay, and compact resolvent in a diffuse algebra gives explicit nonuniformity witnesses. Semicircular families satisfy that special criterion; the general finite-Fisher and rigidity claims retain their missing hypotheses.
\subsection{Sources}
\begin{itemize}
\item Local UnsolvedMath registry, permanent ID \(20002669\), catalogue entry, \texttt{aim-probability-notes.json}, \texttt{notes[110]}, 2006 \emph{Free Analysis} workshop, page 3, section 0.3 ``Free von Neumann Algebras'', question 9: \url{https://aimath.org/WWN/freeanalysis/freeanalysis.pdf}.
\item Supplied research summary identifies the \(L^2\) uniformity through a later question in the same source and Peterson's rigidity framework, and names Dabrowski--Ioana criteria without full bibliographic details or hypotheses.
\item Supplied spectral criterion, compact-resolvent obstruction and semicircular application: attempt 1, partial\_result, partially\_solved, rigor ``proved'', low novelty confidence. The machine-generated aid is not peer reviewed and must be independently verified before citation; no PDF visit or new literature assessment is claimed.
\end{itemize}
\end{document}
